\title{Group Sequence Policy Optimization}

\begin{abstract}
We present Group Sequence Policy Optimization (GSPO), a robust, efficient, and high-performing reinforcement learning algorithm specifically designed for large language model training. Departing from earlier approaches that utilize token-level importance ratios, GSPO instead computes this ratio at the sequence level, enabling sequence-based clipping, reward allocation, and optimization. Through empirical evaluation, we show that GSPO not only surpasses GRPO in both training effectiveness and overall performance but also provides pronounced stability when applied to Mixture-of-Experts (MoE) reinforcement learning. Furthermore, GSPO can potentially streamline the development of RL infrastructure. These advantages of GSPO have played a significant role in the substantial advancements observed in the latest Qwen3 model series.
\end{abstract}

\section{Introduction}
\label{sec:intro}

Reinforcement learning (RL) has become a fundamental approach for advancing the capabilities of language models \citep{o1,dpsk-r1,qwq32b,qwen3}. Employing RL at scale enables language models to address complex challenges—including tasks in advanced mathematics and coding competitions—by engaging in extended and more intricate reasoning.

Achieving scalable reinforcement learning through increased computational resources hinges critically on preserving stable and resilient training processes. Nevertheless, leading RL methods such as GRPO \citep{grpo} encounter pronounced stability challenges during the training of large-scale language models, frequently leading to catastrophic and sometimes irrecoverable model collapse \citep{qwen3, minimax-m1}. These stability limitations present substantial barriers to advancing the capabilities of language models by means of extended RL training.

In this work, we demonstrate that the core issue underlying GRPO's instability is the improper use and subsequent breakdown of importance sampling weights within its algorithmic framework. This flaw generates substantial variance in the training signal, which accumulates as response sequences grow longer and is exacerbated by the clipping strategy, eventually leading to a collapse of the model.

In response to these fundamental constraints, we introduce \textbf{Group Sequence Policy Optimization (GSPO)}, an RL methodology tailored for the training of large language models. The main contribution of GSPO is a rigorously established importance ratio derived from sequence likelihood \citep{click}, which is consistent with the essential framework of importance sampling. Furthermore, GSPO determines normalized rewards by treating them as the advantages across various responses to the same query, thereby guaranteeing consistency between rewards at the sequence level and the underlying optimization process.

Through empirical analysis, we show that GSPO markedly surpasses GRPO in terms of stability during training, efficiency, and achieved performance. Notably, GSPO intrinsically addresses the stability issues typically encountered in reinforcement learning when working with large Mixture-of-Experts (MoE) models, thereby removing the reliance on intricate stabilization methods and indicating the prospect of streamlining RL system architectures. These advantages of GSPO have played a pivotal role in driving the impressive gains realized in the most recent Qwen3 models. We anticipate that GSPO will serve as a solid and scalable algorithmic underpinning for future progress in large-scale reinforcement learning involving language models.

\section{Preliminaries}

\paragraph{Notation}
Throughout this work, we represent an autoregressive language model with parameters $\theta$ as a policy $\pi_\theta$. Queries are symbolized by $x$, and the collection of all queries is indicated by $\mathcal{D}$. For a response $y$ corresponding to a query $x$, its probability under the policy $\pi_\theta$ is expressed as $\pi_\theta (y | x)=\prod_{t=1}^{|y|} \pi_\theta (y_t | x, y_{<t} )$, where $|y|$ specifies the token count in $y$. The combination $(x, y)$, representing a query and its response, can be evaluated by a verifier $r$ to yield a reward $r(x, y) \in [ 0, 1 ]$.

\paragraph{Proximal Policy Optimization (PPO)} Proximal Policy Optimization (PPO) \citep{ppo} utilizes trajectories collected under the previous policy $\pi_{\theta_\text{old}}$ and enforces that policy changes remain within a close neighborhood of this old policy by introducing a clipping strategy. More precisely, PPO optimizes the policy by maximizing the following objective function (here we leave out the KL regularization term for clarity and because it is not central to our discussion):

\begin{align}
\mathcal{J}_\text{PPO}(\theta) = \mathbb{E}_{ x \sim \mathcal{D},\, y \sim \pi_{\theta_\text{old}}( \cdot | x) }
\left[ \frac{1}{|y|} \sum_{t=1}^{|y|}
\min \left( w_{t}(\theta) \widehat{A}_{t},  \, \mathrm{clip} \left( w_{t}(\theta), 1 - {\varepsilon}, 1 + {\varepsilon}\right) \widehat{A}_{t} \right)
\right],
\end{align}

The Proximal Policy Optimization objective is expressed as an expectation with respect to both $x$ sampled from $\mathcal{D}$ and $y$ sampled from the previous policy $\pi_{\theta_\text{old}}( \cdot | x)$. For each sequence of length $|y|$, the objective averages over all time steps $t$. At each step, it considers the minimum between the product $w_{t}(\theta) \widehat{A}_{t}$ and a clipped version, where $w_{t}(\theta)$ is limited to lie within the interval $[1 - \varepsilon, 1 + \varepsilon]$, multiplied by $\widehat{A}_{t}$.

Here, the token $y_t$'s importance ratio is given by $
w_{t}(\theta) = \frac{ \pi_{\theta} (y_{t} | x, y_{<t}) }{ \pi_{\theta_\text{old}} (y_{t} | x,y_{<t})} $, $\widehat{A}_{t}$ denotes the estimated advantage of $y_t$ computed via a separate value model, and $\varepsilon$ represents the allowable clipping threshold for the importance ratios.

A primary practical limitation of PPO stems from its strong dependence on the value model. In particular, the value model is typically as large as the policy model, resulting in substantial memory usage and computational overhead. Moreover, the overall performance of the algorithm is closely tied to the accuracy of the value estimates it provides. Developing a value model that yields reliable estimates is already a difficult problem, and achieving robustness and scalability of these estimates for lengthier outputs and increasingly intricate tasks further compounds this difficulty.

\paragraph{Group Relative Policy Optimization (GRPO)} Rather than relying on a value function, GRPO \citep{grpo} calculates the relative advantage among a set of candidate responses corresponding to a single prompt. In particular, GRPO seeks to optimize the following objective:

\begin{align}
\label{equ:grpo}
\mathcal{J}_\text{GRPO}(\theta) = \mathbb{E}_{ x \sim \mathcal{D},\, \{y_i\}_{i=1}^G \sim \pi_{\theta_\text{old}}( \cdot | x) }
\left[ \frac{1}{G} \sum_{i=1}^{G} \frac{1}{|y_i|} \sum_{t=1}^{|y_i|} 
\min \left( w_{i,t}(\theta) \widehat{A}_{i,t}, \, \mathrm{clip} \left( w_{i,t}(\theta), 1 - \varepsilon, 1 + \varepsilon\right) \widehat{A}_{i,t} \right)
\right],
\end{align}

This objective function for GRPO is formulated as an expectation over input samples $x$ drawn from distribution $\mathcal{D}$ and groups of outputs $\{y_i\}_{i=1}^G$ generated under the policy $\pi_{\theta_\text{old}}(\cdot | x)$. For each group, the average is computed by first summing over all steps $t$ within a group output $y_i$, then normalizing by the output’s length $|y_i|$ and subsequently averaging across all $G$ groups. The inner minimization selects between the value $w_{i,t}(\theta) \widehat{A}_{i,t}$ and its clipped counterpart, where $w_{i,t}(\theta)$ is limited to lie within $[1 - \varepsilon, 1 + \varepsilon]$.

where $G$ is the number of generated responses to each query $x$ (i.e., the group size), and the importance ratio $w_{i,t}(\theta)$ and advantage $\widehat{A}_{i,t}$ of token $y_{i,t}$ are:

\begin{align}
    w_{i,t}(\theta)=\frac{ \pi_{\theta} (y_{i,t} | x, y_{i,<t}) }{ \pi_{\theta_\text{old}} (y_{i,t} | x,y_{i,<t})},\quad\ 
    \widehat{A}_{i,t} = \widehat{A}_{i} = \frac{r(x, y_i) - \operatorname{mean} \left( \{ r(x, y_i) \}_{i=1}^G \right) }{ \operatorname{std} \left( \{ r(x, y_i) \}_{i=1}^G \right) },
\end{align}

respectively, where all the tokens in $y_i$ share the same advantage as $\widehat{A}_{i}$.

\section{Motivation}
\label{sec:motivation}

As model architectures become larger, incorporate greater sparsity (such as in Mixture-of-Experts frameworks), and produce longer outputs, there is an increasing requirement for substantial rollout batch sizes to ensure efficient hardware performance during reinforcement learning. To enhance sample efficiency, it is common to subdivide these sizable rollout batches into several mini-batches, which are then used to compute gradient updates. However, this strategy naturally leads to an off-policy learning scenario: the generated responses $y$ originate from a prior policy $\pi_{\theta_\text{old}}$, not the currently updated policy $\pi_\theta$. As a result, algorithms like PPO and GRPO incorporate clipping mechanisms to mitigate the inclusion of data that diverge too much from the current policy during gradient estimation.

Although techniques such as clipping are designed to address off-policy mismatch, our analysis uncovers a deeper flaw in GRPO: \textit{the objective itself is fundamentally ill-posed}. This issue is magnified during the training of large-scale models on tasks requiring lengthy responses, often culminating in severe model collapse. The root cause of the ill-posedness in the GRPO objective lies in the improper usage of importance sampling weights. Importance sampling is intended to compute the expectation of a function $f$ with respect to a target distribution $\pi_\text{tar}$ by appropriately re-weighting observations obtained from a behavior distribution $\pi_\text{beh}$:

\begin{align}
\mathbb{E}_{z \sim \pi_\text{tar}} \left[ f(z) \right] = \mathbb{E}_{z \sim \pi_\text{beh}} \left[ \frac{ \pi_\text{tar} (z) }{ \pi_\text{beh} (z)} \, f(z) \right].
\end{align}

A key aspect of this approach is the use of a large number of samples ($N \gg 1$) drawn from the behavior distribution $\pi_\text{beh}$. This extensive sampling is necessary for the importance weight $\frac{ \pi_\text{tar} (z) }{ \pi_\text{beh} (z)}$ to successfully mitigate the disparity between the target and behavior distributions.

On the other hand, GRPO incorporates the importance weight $\frac{ \pi_{\theta} (y_{i,t} | x, y_{i,<t}) }{ \pi_{\theta_\text{old}} (y_{i,t} | x,y_{i,<t})}$ at every token step $t$. Because this weighting relies solely on a single sample $y_{i,t}$ drawn from the next-token distribution $\pi_{\theta_\text{old}}( \cdot | x, y_{i, <t})$, it does not serve its intended function of correcting the distribution. Instead, this approach injects significant variance into the gradient estimates during training. Such variance tends to accumulate along extended sequences and is intensified by the clipping procedure. Our experiments reveal that these dynamics can trigger an irreversible form of model collapse. Following such collapse, attempts to recover by continuing training—even when restoring previous checkpoints or exhaustively adjusting hyperparameters like clipping thresholds, increasing sequence length, or modifying the RL setup—do not succeed in rehabilitating the model.

This observation highlights a key limitation inherent in the design of GRPO. The ineffectiveness of token-level importance weights underscores a foundational concept: \textbf{the optimization target must align with the granularity of the reward}. Given that rewards are assigned to full sequences, adjusting for off-policy discrepancies at the token level is inherently flawed. Consequently, we are motivated to discard the token-level objective and instead investigate leveraging importance weights and conducting optimization at the \textit{sequence level} directly.

\section{Algorithm}

\subsection{GSPO: Group Sequence Policy Optimization}

Although GRPO encounters difficulties with the token-level importance weight $\frac{ \pi_{\theta} (y_{i,t} | x, y_{i,<t}) }{ \pi_{\theta_\text{old}} (y_{i,t} | x,y_{i,<t})}$, our findings indicate that within language generation tasks, the \textit{sequence-level} importance weight $\frac{ \pi_{\theta} (y | x) }{ \pi_{\theta_\text{old}} (y | x)}$ possesses a well-defined theoretical interpretation. Specifically, it quantifies the extent to which a response $y$ drawn from $\pi_{\theta_\text{old}} (\cdot | x)$ diverges from the distribution $\pi_{\theta} (\cdot | x)$. This metric is naturally aligned with rewards defined at the sequence level and additionally offers an intuitive measure for understanding the behavior of the clipping strategy.

Based on this straightforward observation, we propose the \textbf{Group Sequence Policy Optimization (GSPO)} algorithm. GSPO employs the following sequence-level optimization objective:

\begin{align}
\mathcal{J}_\text{GSPO} (\theta) = 
\mathbb{E}_{ x \sim \mathcal{D},\, \{y_i\}_{i=1}^G \sim \pi_{\theta_\text{old}}( \cdot | x) }
\left[
\frac{1}{G} \sum_{i=1}^{G}
\min \left( s_{i}(\theta)\,\widehat{A}_{i},\, \mathrm{clip}\left( s_{i}(\theta), 1-{\varepsilon}, 1+{\varepsilon} \right)\,\widehat{A}_{i} \right)
\right],
\label{equ:gspo}
\end{align}

This equation defines the objective function $\mathcal{J}_\text{GSPO} (\theta)$, computed as an expectation over data $x$ drawn from the dataset $\mathcal{D}$ and sequences $\{y_i\}_{i=1}^G$ sampled from the policy $\pi_{\theta_\text{old}}(\cdot|x)$. Within the expectation, the function averages the minimum between $s_{i}(\theta) \widehat{A}_{i}$ and the clipped value $\mathrm{clip}( s_{i}(\theta), 1-{\varepsilon}, 1+{\varepsilon}) \widehat{A}_{i}$ across $G$ samples.

where we adopt the group-based advantage estimation:

\begin{align}
\widehat{A}_{i} = \frac{r(x, y_i) - \mathrm{mean} \left( \{ r(x, y_i) \}_{i=1}^G \right) }{ \mathrm{std} \left( \{ r(x, y_i) \}_{i=1}^G \right) },
\end{align}

and define the importance ratio $s_{i}(\theta)$ based on sequence likelihood \citep{click}:

\begin{align}
s_{i}(\theta) = \left( \frac{ \pi_{\theta} (y_i | x) }{ \pi_{\theta_\text{old}} (y_i | x)} \right)^{\frac{1}{|y_i|}}
=
\exp \left( \frac{1}{|y_i|} \sum_{t=1}^{|y_i|} \log \frac{ \pi_{\theta} (y_{i,t} | x, y_{i,<t}) }{ \pi_{\theta_\text{old}} (y_{i,t} | x,y_{i,<t})} \right).
\end{align}

For each $i$, the score $s_{i}(\theta)$ is computed by taking the ratio of the probabilities assigned to $y_i$ under the new parameterization $\theta$ and the old parameters $\theta_\text{old}$, then applying normalization via the length $|y_i|$. This can equivalently be written as an exponential of the mean, over all time steps $t$ in $y_i$, of the logarithmic difference between the conditional probabilities for each token $y_{i,t}$, conditioned on the history $y_{i,<t}$ and input $x$, under $\theta$ and $\theta_\text{old}$, respectively.

Accordingly, GSPO performs clipping at the response level rather than at the level of individual tokens, in order to eliminate highly "off-policy" samples from the gradient calculation. This approach aligns with both the sequence-level reward evaluation and the corresponding optimization. In addition, length normalization is employed in $s_{i}(\theta)$ so as to mitigate variance and constrain $s_{i}(\theta)$ within a consistent numerical interval. Without this adjustment, small changes in the likelihoods of certain tokens could cause significant swings in the overall sequence-level importance ratio, and as a consequence, the appropriate clipping range would vary depending on response length. Furthermore, it is worth emphasizing that the scale of the clipping range used in GSPO typically diverges by several orders of magnitude from that in earlier methods (such as GRPO), as a direct result of the differing ways in which the importance ratios are formulated.

\subsection{Gradient Analysis}
\label{subsec:gradient}

We can derive the gradient of the GSPO objective as follows (clipping is omitted for brevity):

\begin{align}
\nabla_{\theta} \mathcal{J}_\text{GSPO} (\theta)
=&\
\nabla_{\theta} \mathbb{E}_{ x \sim \mathcal{D},\, \{y_i\}_{i=1}^G \sim \pi_{\theta_\text{old}}( \cdot | x) }
\left[
\frac{1}{G} \sum_{i=1}^{G} s_{i}(\theta) \widehat{A}_{i}
\right] \\
= &\
\mathbb{E}_{ x \sim \mathcal{D},\, \{y_i\}_{i=1}^G \sim \pi_{\theta_\text{old}}( \cdot | x) }
\left[
\frac{1}{G} \sum_{i=1}^{G}
s_{i}(\theta) \widehat{A}_{i}
\cdot \nabla_{\theta} \log s_{i}(\theta)
\right] \\
= &\
\mathbb{E}_{ x \sim \mathcal{D},\, \{y_i\}_{i=1}^G \sim \pi_{\theta_\text{old}}( \cdot | x) }
\left[
\frac{1}{G} \sum_{i=1}^{G}
\left( \frac{ \pi_{\theta} (y_i | x) }{ \pi_{\theta_\text{old}} (y_i | x)} \right)^{\frac{1}{|y_i|}} \widehat{A}_{i}
\cdot \frac{1}{|y_i|} \sum_{t=1}^{|y_i|} \nabla_{\theta} \log \pi_{\theta} (y_{i,t} | x, y_{i,<t})
\right].
\label{equ:gspo_grad}
\end{align}

Paraphrased:

\begin{align}
\nabla_{\theta} \mathcal{J}_\text{GSPO} (\theta)
=&\
\nabla_{\theta} \mathbb{E}_{ x \sim \mathcal{D},\, \{y_i\}_{i=1}^G \sim \pi_{\theta_\text{old}}( \cdot | x) }
\left[
\frac{1}{G} \sum_{i=1}^{G} s_{i}(\theta) \widehat{A}_{i}
\right] \\
= &\
\mathbb{E}_{ x \sim \mathcal{D},\, \{y_i\}_{i=1}^G \sim \pi_{\theta_\text{old}}( \cdot | x) }
\left[
\frac{1}{G} \sum_{i=1}^{G}
s_{i}(\theta) \widehat{A}_{i}
\nabla_{\theta} \log s_{i}(\theta)
\right] \\
=&\
\mathbb{E}_{ x \sim \mathcal{D},\, \{y_i\}_{i=1}^G \sim \pi_{\theta_\text{old}}( \cdot | x) }
\left[
\frac{1}{G} \sum_{i=1}^{G}
\left( \frac{ \pi_{\theta} (y_i | x) }{ \pi_{\theta_\text{old}} (y_i | x)} \right)^{\frac{1}{|y_i|}} \widehat{A}_{i}
\frac{1}{|y_i|} \sum_{t=1}^{|y_i|} \nabla_{\theta} \log \pi_{\theta} (y_{i,t} | x, y_{i,<t})
\right].
\label{equ:gspo_grad}
\end{align}

For comparison, the gradient of the GRPO objective is as follows (note that $\widehat{A}_{i,t} = \widehat{A}_{i}$):

\begin{align}
\nabla_{\theta} \mathcal{J}_\text{GRPO} (\theta) 
=&\ 
\nabla_{\theta} \mathbb{E}_{ x \sim \mathcal{D},\, \{y_i\}_{i=1}^G \sim \pi_{\theta_\text{old}}( \cdot | x) }
\left[ \frac{1}{G} \sum_{i=1}^{G} \frac{1}{|y_i|} \sum_{t=1}^{|y_i|} 
w_{i,t}(\theta) \widehat{A}_{i,t}
\right] \\
=&\
\mathbb{E}_{ x \sim \mathcal{D},\, \{y_i\}_{i=1}^G \sim \pi_{\theta_\text{old}}( \cdot | x) }
\left[ \frac{1}{G} \sum_{i=1}^{G} \widehat{A}_{i} 
\cdot \frac{1}{|y_i|} \sum_{t=1}^{|y_i|} 
\frac{ \pi_{\theta} (y_{i,t} | x, y_{i,<t}) }{ \pi_{\theta_\text{old}} (y_{i,t} | x,y_{i,<t})} 
\nabla_{\theta} \log \pi_{\theta} (y_{i,t} | x, y_{i,<t})  
\right].
\end{align}

Thus, the primary difference between GSPO and GRPO concerns \textit{the method by which the gradients of token log likelihoods are assigned weights}. In the case of GRPO, token gradients are modulated by an ``importance weight'' given by $\frac{ \pi_{\theta} (y_{i,t} | x, y_{i,<t}) }{ \pi_{\theta_\text{old}} (y_{i,t} | x,y_{i,<t})}$. These weights, which may range over $(0, 1+\varepsilon]$ when $\widehat{A}_{i} > 0$ and $[1-\varepsilon, +\infty)$ when $\widehat{A}_{i} < 0$, can differ significantly and are not trivial; as training continues, these disparities can accumulate, introducing instability and potentially unpredictable effects. Conversely, GSPO applies an equal weighting scheme to all tokens within a response, thereby removing this particular source of instability inherent to GRPO.

\subsection{GSPO-token: A Token-level Objective Variant}

In contexts such as multi-turn RL, there may be a need for more detailed advantage modulation at a level finer than the full sequence. To address this, we propose a token-level formulation of GSPO, referred to as \textbf{GSPO-token}, which facilitates the assignment of advantages on a per-token basis:

\begin{align}
\mathcal{J}_\text{GSPO-token}(\theta)
=
\mathbb{E}_{ x \sim \mathcal{D},\, \{y_i\}_{i=1}^G \sim \pi_{\theta_\text{old}}( \cdot | x) }
\bigg[
\frac{1}{G} \sum_{i=1}^{G} \frac{1}{|y_i|} \sum_{t=1}^{|y_i|} 
\min \left( s_{i,t}(\theta) \widehat{A}_{i,t},  \, \mathrm{clip} \left( s_{i,t}(\theta), 1 - {\varepsilon}, 1 + {\varepsilon} \right) \widehat{A}_{i,t} \right) 
\bigg],
\label{equ:gspo-token}
\end{align}

where

\begin{align}
s_{i,t}(\theta) 
= \mathrm{sg} \left[ s_{i}(\theta) \right]  \cdot \frac{ \pi_{\theta} (y_{i,t} | x, y_{i,<t}) }{ \mathrm{sg} \left[ \pi_{\theta} (y_{i,t} | x, y_{i,<t}) \right] },
\end{align}

and $\mathrm{sg}[\cdot]$ denotes only taking the numerical value but stopping the gradient, corresponding to the \texttt{detach} operation in PyTorch. The gradient of GSPO-token can be derived as:

\begin{align}
\nabla_{\theta} \mathcal{J}_\text{GSPO-token}(\theta)
=&\ 
\nabla_{\theta} \mathbb{E}_{ x \sim \mathcal{D},\, \{y_i\}_{i=1}^G \sim \pi_{\theta_\text{old}}( \cdot | x) }
\Bigg[ 
\frac{1}{G} \sum_{i=1}^{G}
\frac{1}{|y_i|} \sum_{t=1}^{|y_i|} 
s_{i,t}(\theta) \widehat{A}_{i,t}
\Bigg] \\
=&\ 
\mathbb{E}_{ x \sim \mathcal{D},\, \{y_i\}_{i=1}^G \sim \pi_{\theta_\text{old}}( \cdot | x) }
\Bigg[ 
\frac{1}{G} \sum_{i=1}^{G} s_i (\theta) \cdot \frac{1}{|y_i|} \sum_{t=1}^{|y_i|} 
\widehat{A}_{i,t} \frac{ \nabla_{\theta} \pi_{\theta} (y_{i,t} | x, y_{i,<t}) }{ \pi_{\theta} (y_{i,t} | x, y_{i,<t}) }
\Bigg] \\
=&\ 
\mathbb{E}_{ x \sim \mathcal{D},\, \{y_i\}_{i=1}^G \sim \pi_{\theta_\text{old}}( \cdot | x) }
\Bigg[ 
\frac{1}{G} \sum_{i=1}^{G}  \left( \frac{ \pi_{\theta} (y_i | x) }{ \pi_{\theta_\text{old}} (y_i | x)} \right)^{\frac{1}{|y_i|}}
\cdot \frac{1}{|y_i|} \sum_{t=1}^{|y_i|}
\widehat{A}_{i,t} \nabla_{\theta} \log \pi_{\theta} (y_{i,t} | x, y_{i,<t})  
\Bigg].
\label{equ:gspo-token_grad}
\end{align}

Observe that the expression $\frac{ \pi_{\theta} (y_{i,t} | x, y_{i,<t}) }{ \mathrm{sg} \left[ \pi_{\theta} (y_{i,t} | x, y_{i,<t}) \right] }$ evaluates to 1 for any input; consequently, $s_{i,t}(\theta)$ is, in practice, equivalent in value to $s_{i}(\theta)$.

If one contrasts Equation~\eqref{equ:gspo} with Equation~\eqref{equ:gspo-token}, as well as Equation~\eqref{equ:gspo_grad} with Equation~\eqref{equ:gspo-token_grad}, it becomes apparent that both GSPO and GSPO-token yield identical results for the optimization objective, clipping constraints, and analytic gradients whenever the advantages $\widehat{A}_{i,t}$ for each token in $y_i$ are set uniformly equal to $\widehat{A}_{i}$. However, the GSPO-token formulation offers enhanced adaptability, as it permits specification of token-level advantage values independently.

\section{Experiments and Discussion}

\subsection{Empirical Results}

We utilize a cold-start setup where the model is fine-tuned from Qwen3-30B-A3B-Base, presenting both the training reward trajectories and evaluation results on three benchmarks: AIME'24 (Pass@1 averaged over 32 generations), LiveCodeBench (202410-202502, Pass@1 averaged across 8 samples), and CodeForces (Elo Rating). Throughout RL training, each rollout batch is subdivided into four mini-batches for the purpose of gradient computation. For GSPO, we configure the left and right clipping bounds in Equation~\eqref{equ:gspo} to 3e-4 and 4e-4, respectively. As a point of comparison, we evaluate against GRPO, which serves as our baseline; its left and right clipping bounds in Equation~\eqref{equ:grpo} are set to 0.2 and 0.27, as determined by careful parameter tuning to maintain experimental fairness. Importantly, GRPO requires the Routing Replay training mechanism for successful MoE RL convergence, which we further elaborate upon in \S~\ref{subsec:routing_replay}; by contrast, \textit{GSPO eliminates the dependency on this approach}.

Figure~\ref{fig:results} illustrates that the GSPO training process maintains consistent stability throughout. Our findings indicate that \textbf{GSPO enables ongoing gains in performance by scaling up training compute, systematically refreshing the query set, and lengthening the generation outputs}. In addition, GSPO exhibits greater training efficiency compared to GRPO, yielding higher training accuracy and superior benchmark results when utilizing equivalent amounts of training computation and query consumption. Lastly, we have effectively integrated GSPO into the RL training pipeline for the most recent Qwen3 models, offering robust evidence that GSPO is instrumental in fully leveraging RL scaling for large language models.

\subsection{Curious Observation on Clipping Fractions}

A primary difference between GSPO and GRPO lies in the fact that GSPO implements clipping at the response level, rather than at the level of individual tokens. Specifically, as illustrated in Figure~\ref{fig:clipping}, the proportion of tokens clipped by GSPO and GRPO differs by approximately two orders of magnitude, a gap that persists regardless of any adjustments to the clipping ranges. Nonetheless, GSPO—despite discarding a significantly higher number of tokens from the training process (either in terms of training or gradient calculation)—demonstrates greater training efficiency relative to GRPO. This unexpected outcome, in which a higher rate of token clipping results in improved training efficiency, suggests that the token-level gradient estimates provided by GRPO are fundamentally noisy and less effective for making use of samples. Conversely, the sequence-level methodology adopted in GSPO yields a more consistent and robust learning signal.

\subsection{Benefit of GSPO for MoE Training}
\label{subsec:routing_replay}

\paragraph{Background}
Unlike RL training in dense models, MoE models, with their sparsely activated experts, face distinctive challenges related to training stability. Specifically, when employing the GRPO algorithm, the inherently variable \textit{expert-activation volatility} observed in MoE models disrupts the convergence of RL training. Concretely, the set of experts invoked for producing the same model output may shift considerably after each gradient update. For instance, in the 48-layer Qwen3-30B-A3B-Base model, analysis reveals that about 10\% of the experts selected for the identical rollout sample change following each RL gradient update, comparing the experts activated by the current policy $\pi_{\theta}$ and those by the previous policy $\pi_{\theta_\text{old}}$.
This issue, which intensifies as the depth of MoE models increases, leads to large variations in the token-level importance ratios $w_{i,t}(\theta) = \frac{ \pi_{\theta} (y_{i,t} | x, y_{i,<t}) }{ \pi_{\theta_\text{old}} (y_{i,t} | x,y_{i,<t})}$, severely diminishing their reliability (refer to \S~\ref{sec:motivation} and~\ref{subsec:gradient}) and impeding stable RL training convergence.

\paragraph{Our Previous Approach} In our earlier work addressing this problem, we utilized the \textbf{Routing Replay} training technique. In this method, the experts activated according to $\pi_{\theta_\text{old}}$ are cached, and these same routing decisions are later ``replayed'' within $\pi_{\theta}$ when evaluating the importance weights $w_{i,t}(\theta) = \frac{ \pi_{\theta} (y_{i,t} | x, y_{i,<t}) }{ \pi_{\theta_\text{old}} (y_{i,t} | x,y_{i,<t})}$. This ensures that for any token $y_{i,t}$, both $\pi_{\theta} (y_{i,t} | x, y_{i,<t})$ and $\pi_{\theta_\text{old}} (y_{i,t} | x,y_{i,<t})$ operate under the identical active subnetwork, thus stabilizing token-level importance ratios and preserving consistent optimization of the underlying activated network throughout the course of gradient steps. As illustrated in Figure~\ref{fig:routing_replay}, Routing Replay proves vital for achieving proper convergence during GRPO training in MoE architectures.

\paragraph{Benefit of GSPO} While Routing Replay supports GRPO training of MoE models to achieve convergence, this approach of recycling routing patterns introduces significant memory and communication costs, and can also restrict the effective expressivity of the MoE model. In comparison, as depicted in Figure~\ref{fig:results}, GSPO removes reliance on Routing Replay, enabling conventional computation of importance ratios $s_i(\theta)$ while maintaining robust convergence and steady optimization. The primary reason is that GSPO targets the sequence-level likelihood $\pi_{\theta} (y_i | x)$ and remains largely unaffected by changes in the token-level likelihood $\pi_{\theta} (y_{i,t} | x, y_{i,<t})$. Since MoE models consistently retain their language modeling proficiency, fluctuations in sequence likelihood are minimal. Therefore, GSPO directly addresses the issue of volatile expert activation inherent to MoE models, eliminating the need for intricate solutions such as Routing Replay. This advancement both streamlines and stabilizes training, and empowers the model to fully exploit its capacity, unencumbered by artificial limits.

\subsection{Benefit of GSPO for RL Infrastructure}

Due to differences in precision between training engines (such as Megatron) and inference engines (like SGLang and vLLM), it is common practice to recalculate the likelihoods of generated responses under the previous policy $\pi_{\theta_\text{old}}$ using the training engine. In contrast, GSPO bases its optimization on sequence-level likelihoods rather than token-level ones, and intuitively, sequence-level measures are far less sensitive to these precision variations. As a result, GSPO allows the direct use of likelihood outputs from the inference engine for optimization purposes, removing the necessity of recomputation with the training engine. This advantage is particularly useful for scenarios involving partial rollout, multi-turn reinforcement learning, and architectures in which training and inference are decoupled.

\section{Conclusion}

We introduce Group Sequence Policy Optimization (GSPO), a novel reinforcement learning method designed for training large language models. Building on the foundation of importance sampling, GSPO establishes importance ratios grounded in sequence likelihood and applies clipping, reward assignment, and optimization at the sequence level. In empirical evaluations, GSPO delivers markedly improved training robustness, higher efficiency, and enhanced outcomes over GRPO, and proves especially effective when applied to large-scale reinforcement learning of MoE models—contributing to the significant advancements realized in recent Qwen3 models. As a scalable and reliable algorithmic platform, GSPO enables further scaling of reinforcement learning, which we anticipate will catalyze major progress in the advancement of artificial intelligence.

\end{document}